% Bagian: turing-machines
% Bab: undecidability
% Subbagian: trakhtenbrot

\documentclass[../../../include/open-logic-section]{subfiles}

\begin{document}

\olfileid[id]{tur}{und}{tra}
\olsection{Teorema Trakhtenbrot}

\begin{explain}
  Dalam \olref[rep]{sec}, kita mendefinisikan !!{sentence}
  $!T(M,w)$ dan~$!E(M,w)$ untuk mesin Turing~$M$ dan untai masukan~$w$.
  Kemudian dalam \olref[ver]{lem:valid-if-halt} dan
  \olref[ver]{lem:halt-if-valid}, kita menunjukkan bahwa
  $!T(M,w) \lif !E(M,w)$ valid jika dan hanya jika $M$, ketika dijalankan
  pada masukan~$w$, akhirnya berhenti. Karena Masalah Penghentian tidak dapat
  diputuskan, ini menyiratkan bahwa validitas dan keterpenuhan !!{sentence}
  logika orde pertama tidak dapat diputuskan
  (\cref{tur:und:uns:thm:decision-prob,tur:und:uns:cor:undecidable-sat}).

  Akan tetapi, validitas dan keterpenuhan kalimat didefinisikan terhadap
  !!{structure} sebarang, baik berhingga maupun tak berhingga. Anda mungkin
  menduga bahwa lebih mudah memutuskan apakah !!a{sentence} dapat dipenuhi
  dalam !!a{structure} berhingga (atau valid dalam semua !!{structure}
  berhingga). Kita dapat mengadaptasi bukti ketakterputusan masalah keputusan
  untuk menunjukkan bahwa dugaan itu tidak benar.

  Pertama, jika Anda kembali ke bukti \olref[ver]{lem:halt-if-valid}, Anda akan
  melihat bahwa di sana kita menghasilkan model~$\Struct M$ dari~$!T(M,w)$
  yang mendeskripsikan secara tepat apa yang dilakukan mesin~$M$ ketika
  dijalankan pada masukan~$w$. Domain model itu adalah~$\Nat$, jadi tak
  berhingga. Namun, jika $M$ benar-benar berhenti pada masukan~$w$, kita dapat
  membangun model berhingga~$\Struct M'$ dengan cara serupa. Andaikan $M$ yang
  dijalankan pada masukan~$w$ berhenti setelah~$k$ langkah. Ambil sebagai
  domain~$\Domain{M'}$ himpunan $\{0, \dots, n\}$, dengan
  $n = \max(k+1,\len{w})$, lalu tetapkan
  \[
    \Assign{\prime}{M'}(x) =
    \begin{cases}
      x + 1 &\text{jika $x < n$}\\
      n &\text{jika tidak,}
    \end{cases}
  \]
  dan $\tuple{x,y} \in \Assign{<}{M'}$ jika dan hanya jika $x < y$ atau
  $x = y = n$. Selain itu, $\Struct{M'}$ didefinisikan sama seperti
  $\Struct{M}$. Berdasarkan definisi~$\Struct{M'}$, sebagaimana dalam bukti
  \olref[ver]{lem:halt-if-valid}, $\Sat{M'}{!T(M,w)}$. Karena kita
  mengandaikan bahwa $M$ berhenti pada masukan~$w$,
  $\Sat{M'}{!E(M,w)}$. Jadi, $\Struct{M'}$ adalah model berhingga dari
  $!T(M,w) \land !E(M,w)$ (perhatikan bahwa kita mengganti $\lif$
  dengan~$\land$).

  Kita sudah menempuh separuh bukti: telah kita tunjukkan bahwa jika $M$
  berhenti pada masukan~$w$, maka $!T(M,w) \land !E(M,w)$ memiliki model
  berhingga. Sayangnya, arah sebaliknya tidak berlaku. Artinya, terdapat mesin
  Turing yang tidak berhenti pada suatu masukan~$w$, tetapi
  $!T(M,w) \land !E(M,w)$ tetap memiliki model berhingga. Sebagai contoh,
  perhatikan mesin~$M$ dengan satu-satunya keadaan $q_0$ dan instruksi
  $\delta(q_0,\TMblank) = \tuple{q_0,\TMblank,\TMstay}$. Ketika dijalankan
  pada masukan kosong~$w = \emptyseq$, mesin ini tidak pernah berhenti: mesin
  berada dalam kalang tak berhingga, tetapi tidak mengubah pita maupun
  menggerakkan kepala. Semua konfigurasinya sama (keadaan, posisi kepala, dan
  isi pita sama). Kita dapat mendefinisikan !!a{structure} berhingga
  $\Struct{M''}$ yang memenuhi $!T(M,\emptyseq) \land !E(M,\emptyseq)$
  (latihan). Akan tetapi, kita dapat mengubah~$!T(M,w)$ dengan cara yang sesuai
  sehingga !!{structure} semacam itu tersingkir.

  \begin{prob}
    Misalkan $M$ adalah mesin Turing dengan satu-satunya keadaan $q_0$ dan
    satu-satunya instruksi
    $\delta(q_0,\TMblank) = \tuple{q_0,\TMblank,\TMstay}$.
    Misalkan $\Domain{M''} = \{0, 1, 2\}$,
    $\Assign{\prime}{M''}(0) = \Assign{\prime}{M''}(1) = 1$ dan
    $\Assign{\prime}{M''}(2) = 2$, serta
    $\Assign{<}{M''} = \{\tuple{0,1}, \tuple{1,1}, \tuple{2,2}\}$.
    Definisikan $\Assign{\Obj Q_{q_0}}{M''}$,
    $\Assign{\Obj S_{\TMblank}}{M''}$, dan
    $\Assign{\Obj S_{\TMendtape}}{M''}$ sedemikian sehingga
    $!T(M,\emptyseq)$ dan $!E(M,\emptyseq)$ menjadi benar, lalu jelaskan
    alasannya. Petunjuk: perhatikan bahwa $\delta(q_0, \TMendtape)$ tidak
    terdefinisi. Pastikan bahwa
    \begin{align*}
    & \Obj Q_{q_0}(\num{1}, \num{n}) \land \Obj S_{\TMendtape}(\num{0}, \num{n})
      \land
      \lforall[x][(\num{0} < x \lif \Obj S_\TMblank(x, \num{n}))]
      \text{\quad untuk semua $n \in \Nat$}\\
    & \lexists[y][(\Obj Q_{q_0}(\num{0}, y) \land \Obj S_{\TMendtape}(\num{0}, y))]
    \end{align*}
    keduanya benar dalam~$\Struct{M''}$.
  \end{prob}
\end{explain}

Perhatikan !!{sentence} yang mendeskripsikan operasi mesin Turing~$M$ pada
masukan $w = \sigma_{i_1}\dots\sigma_{i_k}$:
\begin{enumerate}
\item Aksioma yang mendeskripsikan bilangan dan~$<$ (sama seperti dalam
definisi~$!T(M,w)$ di \olref[rep]{sec}).
\item Aksioma yang mendeskripsikan konfigurasi masukan: sama seperti dalam
definisi~$!T(M,w)$.
\item Aksioma yang mendeskripsikan transisi dari satu konfigurasi ke
  konfigurasi berikutnya:

Untuk uraian berikut, misalkan $!A(x, y)$ seperti sebelumnya, dan misalkan
\[
  !B(y) \ident \lforall[x][(x < y \lif \eq/[x][y])].
\]
\begin{enumerate}
\item \ollabel{rep-right} Untuk setiap instruksi $\delta(q_i, \sigma) =
  \tuple{q_j, \sigma', \TMright}$, !!{sentence} tersebut adalah:
\begin{align*}
& \lforall[x][\lforall[y][(
   (\Obj Q_{q_i}(x, y) \land \Obj S_{\sigma}(x, y)) \lif {}]] \\
&\qquad   (\Obj Q_{q_j}(x', y') \land \Obj S_{\sigma'}(x, y') \land
!A(x, y) \land !B(y')))
\end{align*}
\item \ollabel{rep-left} Untuk setiap instruksi $\delta(q_i, \sigma) =
  \tuple{q_j, \sigma', \TMleft}$, !!{sentence} tersebut adalah
\begin{align*}
& \lforall[x][\lforall[y][
    ((\Obj Q_{q_i}(x', y) \land \Obj S_{\sigma}(x', y)) \lif {}]]\\
& \qquad   (\Obj Q_{q_j}(x, y') \land \Obj S_{\sigma'}(x', y') \land
!A(x', y) \land !B(y'))) \land {}\\
& \lforall[y][((\Obj Q_{q_i}(\num{0}, y) \land \Obj S_{\sigma}(\num{0},
    y)) \lif {}]\\
& \qquad (\Obj Q_{q_j}(\num{0}, y') \land \Obj S_{\sigma'}(\num{0},
  y') \land !A(\num{0}, y) \land !B(y')))
\end{align*}
\item \ollabel{rep-stay} Untuk setiap instruksi $\delta(q_i, \sigma) =
  \tuple{q_j, \sigma', \TMstay}$, !!{sentence} tersebut adalah:
\begin{align*}
& \lforall[x][\lforall[y][(
   (\Obj Q_{q_i}(x, y) \land \Obj S_{\sigma}(x, y)) \lif {}]] \\
&\qquad (\Obj Q_{q_j}(x, y') \land \Obj S_{\sigma'}(x, y') \land
!A(x, y) \land !B(y')))
\end{align*}
\end{enumerate}
Seperti yang dapat Anda lihat, !!{sentence} yang mendeskripsikan transisi~$M$
sama dengan !!{sentence} yang bersesuaian dalam~$!T(M,w)$, kecuali bahwa kita
menambahkan $!B(y')$ pada bagian akhir. $!B(y')$ memastikan bahwa bilangan
$y'$ bagi konfigurasi ``berikutnya'' berbeda dari semua bilangan sebelumnya
$\Obj 0$, $\Obj 0'$, \dots.
% \item Sentences that express consistency conditions:
% \begin{enumerate}
%   \item\ollabel{rep-con-symb}%
%   !!^a{sentence} that says that no square can have more than one
%   symbol on it at any given time:
%   \[\lforall[x][\lforall[y][(\Obj S_{\sigma}(x, y) \lif \lnot \Obj
%   S_{\sigma'}(x, y))]],\]
%   for every pair $\sigma \neq \sigma'$ of tape symbols of~$T$.
%   \item\ollabel{rep-con-state}%
%   !!^a{sentence} that says that $M$ cannot be in more than one state
%   at any given time:
%   \[\lforall[x][\lforall[y][(\Obj Q_{q}(x, y) \lif
%   \lforall[z][\lnot \Obj
%   Q_{q'}(z, y))]],\]
%   for every pair $q \neq q'$ of states of~$T$.
%   \item\ollabel{rep-con-tape}%
%   !!^a{sentence} that says that $M$ cannot be on more than one tape
%   square at any given time:
%   \[\lforall[x][\lforall[y][(\Obj Q_{q}(x, y) \lif
%   \lforall[z][\eq/[z][y]\lif \lnot \Obj
%   Q_{q'}(z, y))]],\]
%   for every pair $q$, $q'$ of states of~$T$.
% \end{enumerate}
\end{enumerate}
Misalkan $!T'(M, w)$ adalah konjungsi semua !!{sentence} di atas bagi mesin
Turing~$M$ dan masukan~$w$.

\begin{lem}\ollabel{lem:halts-sat}
  Jika $M$ yang dijalankan pada masukan~$w$ berhenti, maka
  $!T'(M,w) \land !E(M,w)$ memiliki model berhingga.
\end{lem}

\begin{proof}
  Misalkan $\Struct{M'}$ seperti dalam bukti
  \olref[ver]{lem:halt-if-valid}, kecuali
  \begin{align*}
    \Domain{M'} & = \{0, \dots, n\},\\
    \Assign{\prime}{M'}(x) & =
    \begin{cases}
      x + 1 &\text{jika $x < n$}\\
      n &\text{jika tidak,}
    \end{cases} \\
    \tuple{x,y} \in \Assign{<}{M'} &\text{ jika dan hanya jika
      $x < y$ atau $x = y = n$,}
  \end{align*}
  dengan $n = \max(k+1,\len{w})$ dan $k$ adalah bilangan terkecil sedemikian
  sehingga $M$ yang dijalankan pada masukan~$w$ telah berhenti setelah~$k$
  langkah. Verifikasi bahwa $\Sat{M'}{!T'(M,w) \land !E(M,w)}$ kita serahkan
  sebagai latihan.
\end{proof}

\begin{prob}
  Lengkapilah bukti \olref[tur][und][tra]{lem:halts-sat} dengan membuktikan
  bahwa $\Sat{M'}{!T'(M,w) \land !E(M,w)}$.
\end{prob}

\begin{lem}\ollabel{lem:sat-halts}
  Jika $!T'(M,w) \land !E(M,w)$ memiliki model berhingga, maka $M$ yang
  dijalankan pada masukan~$w$ berhenti.
\end{lem}

\begin{proof}
  Kita menunjukkan kontraposisinya. Andaikan $M$ yang dijalankan pada~$w$
  tidak berhenti. Jika $!T'(M,w) \land !E(M,w)$ sama sekali tidak memiliki
  model, kita selesai. Jadi, andaikan $\Struct{M}$ adalah model dari
  $!T'(M,w) \land !E(M,w)$. Kita harus menunjukkan bahwa model itu tidak dapat
  berhingga.

  Kita dapat membuktikan, sama seperti dalam \olref[ver]{lem:config}, bahwa
  jika~$M$, ketika dijalankan pada masukan~$w$, belum berhenti setelah~$n$
  langkah, maka
  $!T'(M,w) \Entails !C(M, w, n) \land !B(\num{n})$. Karena $M$ yang
  dijalankan pada masukan~$w$ tidak berhenti,
  $!T'(M,w) \Entails !C(M, w, n) \land !B(\num{n})$ untuk semua
  $n \in \Nat$. Perhatikan bahwa menurut \olref[rep]{prop:mlessk},
  $!T'(M,w) \Entails \num{k} < \num{n}$ untuk semua~$k < n$. Selain itu,
  $!B(\num{n}) \Entails \num{k} < \num{n} \lif
  \eq/[\num{k}][\num{n}]$. Jadi,
  $\Sat{M}{\eq/[\num{k}][\num{n}]}$ untuk semua~$k < n$; yakni, tak berhingga
  banyaknya suku~$\num{k}$ harus memiliki nilai yang semuanya berbeda dalam
  $\Struct{M}$. Namun, hal ini mengharuskan $\Domain{M}$ tak berhingga, sehingga
  $\Struct{M}$ tidak dapat menjadi model berhingga dari
  $!T'(M,w) \land !E(M,w)$.
\end{proof}

\begin{prob}
  Lengkapilah bukti \olref[tur][und][tra]{lem:sat-halts} dengan membuktikan
  bahwa jika~$M$, ketika dijalankan pada masukan~$w$, belum berhenti setelah
  $n$ langkah, maka $!T'(M,w) \Entails !B(\num{n})$.
\end{prob}

\begin{thm}[Teorema Trakhtenbrot]
  \ollabel{thm:trakhtenbrodt}
  Tidak dapat diputuskan apakah suatu !!{sentence} sebarang logika orde
  pertama memiliki model berhingga (yakni dapat dipenuhi secara berhingga).
\end{thm}

\begin{proof}
  Andaikan terdapat mesin Turing~$F$ yang memutuskan masalah keterpenuhan
  berhingga. Maka, jika diberikan mesin Turing~$M$ dan masukan~$w$, kita dapat
  menghitung kalimat~$!T'(M,w) \land !E(M,w)$ lalu menggunakan~$F$ untuk
  memutuskan apakah kalimat itu memiliki model berhingga. Menurut
  \cref{tur:und:tra:lem:halts-sat,tur:und:tra:lem:sat-halts}, kalimat itu
  memiliki model berhingga jika dan hanya jika $M$ yang dijalankan pada
  masukan~$w$ berhenti. Jadi, $F$ dapat kita gunakan untuk menyelesaikan
  masalah penghentian, yang kita ketahui tidak dapat diselesaikan.
\end{proof}

\begin{cor}\ollabel{cor:fproof-incomp}%
  Tidak mungkin ada sistem !!{derivation} yang sahih dan lengkap bagi
  validitas \emph{berhingga}, yakni sistem !!a{derivation} yang memenuhi
  $\Proves !B$ jika dan hanya jika $\Sat{M}{!B}$ untuk setiap
  !!{structure} berhingga~$\Struct{M}$.
\end{cor}

\begin{proof}
  Latihan.
\end{proof}

\begin{prob}
  Buktikan \olref[tur][und][tra]{cor:fproof-incomp}. Perhatikan bahwa $!B$
  dipenuhi dalam setiap !!{structure} berhingga jika dan hanya jika
  $\lnot !B$ tidak dapat dipenuhi secara berhingga. Jelaskan mengapa
  keterpenuhan berhingga bersifat semidapat diputuskan dalam pengertian
  \olref[tur][und][uns]{thm:valid-ce}. Gunakan hal ini untuk menunjukkan bahwa
  jika terdapat sistem !!a{derivation} bagi validitas berhingga, maka
  keterpenuhan berhingga dapat diputuskan.
\end{prob}


\end{document}
